\documentclass[11pt]{article}
\usepackage[a4paper,margin=1.9cm]{geometry}
\usepackage[T1]{fontenc}
\usepackage{textcomp}
\usepackage[spanish]{babel}
\usepackage{amsmath,amssymb}
\usepackage{enumitem}
\usepackage{tikz}
\usetikzlibrary{calc,arrows.meta,patterns,decorations.pathmorphing}
\usepackage{xcolor}
\usepackage{multicol}
\usepackage{parskip}
\usepackage{lmodern}
\renewcommand{\familydefault}{\sfdefault}

\definecolor{unalblue}{RGB}{31,73,124}

\newlist{opciones}{enumerate}{1}
\setlist[opciones]{label=(\alph*),itemsep=1pt,topsep=2pt,leftmargin=1.8em,font=\normalfont}
\newcommand{\ptsbox}[1]{\fbox{\footnotesize #1}~}

\newcommand{\vect}[2]{\begin{pmatrix}#1\\#2\end{pmatrix}}
\newcommand{\vectiii}[3]{\begin{pmatrix}#1\\#2\\#3\end{pmatrix}}

\newcommand{\examheader}{%
  \noindent
  \begin{minipage}[t]{0.6\linewidth}
    {\small\bfseries Geometría Vectorial y Analítica --- Parcial 2}\\
    {\small Semestre 2015--01}
  \end{minipage}%
  \hfill
  \begin{minipage}[t]{0.35\linewidth}
    \raggedleft\small Escuela de Matemáticas. Universidad Nacional de Colombia, Sede Medellín.
  \end{minipage}\\[2pt]
  \rule{\linewidth}{0.6pt}
}
\newcommand{\examfooter}{%
  \vfill
  \rule{\linewidth}{0.4pt}\\[1pt]
  {\scriptsize\textit{Transcripción digital --- MathUNAL}\hfill\textcolor{unalblue}{\textbf{\thepage}}}
}
\pagestyle{empty}

\begin{document}

\examheader

\begin{center}
{\bfseries Segundo Examen Parcial de Geometría}\\
18 de Abril del 2015

\vspace{4pt}
\textbf{Puntaje. Sólo para uso Oficial}

\vspace{2pt}
\begin{tabular}{|c|c|c|c|c|c|c|}
\hline
1 (20) & 2 (16) & 3 (23) & 4 (20) & 5 (21) & Total & Nota\\
\hline
 & & & & & & \\
\hline
\end{tabular}
\end{center}

\textbf{Instrucciones:} Antes de empezar verifique que su examen esté completo y que sus datos sean correctos. Por favor verifique que su celular esté apagado. No está permitido el uso de calculadora ni de otros aparatos electrónicos. Solucione las preguntas en los espacios en blanco asignados a cada una de ellas. No está permitido sacar hojas en blanco ni ningún tipo de apuntes durante el examen. Utilice la página 6 para realizar sus operaciones en borrador. Duración: \textbf{1 hora y 50 minutos}.

\vspace{10pt}

\begin{enumerate}[leftmargin=1.6em]

\item Considere la transformación lineal dada por $T\vect{x}{y} = \vect{6x}{x+2y}$.

\begin{opciones}
\item \ptsbox{5\%} Encuentre los valores propios de $T$.

\vspace{50pt}

\item \ptsbox{7\%} Encuentre el conjunto de vectores propios asociado a cada valor propio de $T$.

\vspace{60pt}

\item \ptsbox{8\%} Determine si es posible factorizar la matriz asociada a $T$, $m(T)$, en la forma $m(T) = PDP^{-1}$ con $D$ diagonal y $P$ invertible. Encuentre $P$, $D$ y $P^{-1}$.

\vspace{70pt}
\end{opciones}

\item Para cada una de las transformaciones del plano dadas en los siguientes literales diga si la transformación es un movimiento euclidiano justificando la respuesta. (Nota: no se conceden puntos a respuestas sin justificación).

\begin{opciones}
\item \ptsbox{8\%} $T\vect{x}{y} = \vect{-1}{2} + \vect{x+y}{-x+y}$.

\vspace{60pt}

\item \ptsbox{8\%} $R\vect{x}{y} = \vect{(3/\sqrt{10})x - (1/\sqrt{10})y}{(1/\sqrt{10})x + (3/\sqrt{10})y}$.

\vspace{60pt}
\end{opciones}

\newpage
\examheader

\item \ptsbox{23\%} Sea $\mathcal{T}_1$ el triángulo con vértices $A = \vect{0}{0}$, $B = \vect{1}{3}$ y $C = \vect{3}{1}$.

\begin{opciones}
\item \ptsbox{7\%} Describa paramétricamente el triángulo $\mathcal{T}_1$.

\begin{center}
\begin{tikzpicture}[scale=0.9]
\coordinate (A) at (0,0);
\coordinate (B) at (1,3);
\coordinate (C) at (3,1);
\draw[-{Stealth[length=2mm]}] (-0.5,0) -- (4,0) node[right] {\small $x$};
\draw[-{Stealth[length=2mm]}] (0,-0.5) -- (0,4) node[above] {\small $y$};
\draw[thick] (A) -- (B) -- (C) -- cycle;
\node[below left] at (A) {\small $A$};
\node[above] at (B) {\small $B$};
\node[right] at (C) {\small $C$};
\node at (1.3,1.3) {\small $\mathcal{T}_1$};
\node[left] at (0.5,1.5) {\small $L_1$};
\node[above] at (2,2) {\small $L_2$};
\node[below] at (1.5,0.5) {\small $L_3$};
\end{tikzpicture}
\end{center}

\vspace{20pt}

\item \ptsbox{8\%} Halle el área de $\mathcal{T}_1$.

\vspace{50pt}

\item \ptsbox{8\%} Sea $T$ la transformación lineal cuya matriz es $m(T) = \begin{pmatrix} 2 & -1 \\ 5 & 1 \end{pmatrix}$, y sea $\mathcal{T}_2$ el triángulo tal que $T(\mathcal{T}_2) = \mathcal{T}_1$. Halle el área de $\mathcal{T}_2$.

\vspace{50pt}
\end{opciones}

\item \ptsbox{20\%} Sea $\mathcal{P}_1$ el paralelogramo generado por los vectores $\vect{2}{1}$ y $\vect{-1}{3}$, y sea $\mathcal{P}_2$ el paralelogramo generado por los vectores $\vect{-1}{2}$ y $\vect{3}{1}$. Encuentre la matriz $A$ de una transformación lineal $T$ tal que $T(\mathcal{P}_1) = \mathcal{P}_2$.

\vspace{90pt}

\end{enumerate}

\examfooter
\end{document}
